\section{An explicit quadratic formula}\label{sec:quadratic}
Let $C^{(a)}$ be the positive prolate operator on $L^2(-1,1)$ with
kernel\label{def:ka}
$k_a(x-y)=\sin(a(x-y))/(\pi(x-y))$, and put
$C_\tau=C^{(\pi\tau)}$, $C_{2\tau}=C^{(2\pi\tau)}$.
Write $D_R=\operatorname{tr}(C_{2\tau}-C_{2\tau}^2)$,
$D_J=4\tau-\operatorname{tr}\mathcal K_\tau^2$,
$\gamma$ for Euler's constant, $\Si$ as in \eqref{eq:2.8}, and
$\operatorname{Ci}(x)=-\int_x^\infty\cos(t)/t\,dt$.
The following independent calculation determines a persistent oscillation
that the general bounded-remainder theorem does not specify.
\begin{lemma}[Exact defect trace of the sinc kernel]\label{lem:defect-exact}
For every $a>0$,
\begin{equation}\label{eq:defect-exact}
 \operatorname{tr}\bigl(C^{(a)}-(C^{(a)})^2\bigr)
 =\frac1{\pi^2}\Bigl[\log(4a)+1+\gamma-\operatorname{Ci}(4a)-2\mathfrak J(a)\Bigr],
\end{equation}
where $\mathfrak J(a)=\int_2^\infty\cos(2au)\,u^{-2}\,du
=\tfrac12\cos(4a)-2a\bigl(\tfrac\pi2-\Si(4a)\bigr)$, with $|\operatorname{Ci}(4a)|\le 1/(2a)$ and $|\mathfrak J(a)|\le 1/(4a)$.
Consequently
\begin{equation}\label{eq:defect-asymp}
 \Bigl|\operatorname{tr}\bigl(C^{(a)}-(C^{(a)})^2\bigr)
 -\frac{\log(4a)+1+\gamma}{\pi^2}\Bigr|\le\frac1{\pi^2a},
 \qquad
 D_R=\frac{\log(8\pi\tau)+1+\gamma}{\pi^2}+O(\tau^{-1}).
\end{equation}
\end{lemma}

\begin{proof}
The displacement variable and Plancherel give
$\tr(C^{(a)})^2=\int_{-2}^2(2-|u|)k_a(u)^2\,du$ and
$\tr C^{(a)}=4\int_0^\infty k_a(u)^2\,du=2a/\pi$. Hence
\[
 \tr\bigl(C^{(a)}-(C^{(a)})^2\bigr)
 =4\int_2^\infty k_a(u)^2\,du+2\int_0^2u\,k_a(u)^2\,du
 =\frac{1-2\mathfrak J(a)+\gamma+\log(4a)-\operatorname{Ci}(4a)}{\pi^2},
\]
using $2\sin^2(au)=1-\cos(2au)$ and
$\int_0^x(1-\cos t)t^{-1}\,dt=\gamma+\log x-\operatorname{Ci}(x)$.
Integration by parts in the two directions gives
\[
\begin{split}
 \operatorname{Ci}(x)&=\frac{\sin x}{x}-\int_x^\infty\frac{\sin t}{t^2}\,dt,\\
 \mathfrak J(a)&=-\frac{\sin(4a)}{8a}
       +\frac1a\int_2^\infty\frac{\sin(2au)}{u^3}\,du
 =\frac{\cos(4a)}2-2a\bigl(\tfrac\pi2-\Si(4a)\bigr).
\end{split}
\]
Thus $|\operatorname{Ci}(4a)|\le1/(2a)$ and
$|\mathfrak J(a)|\le1/(4a)$, proving all the stated bounds.
\end{proof}

\begin{remark}\label{rem:defect-remainder}
The leading $1/a$ contributions from $-\operatorname{Ci}(4a)$ and $-2\mathfrak J(a)$
cancel.  Further integration by parts gives the remainder
$-\cos(4a)/(16\pi^2a^2)+O(a^{-3})$ in \eqref{eq:defect-asymp}.
We use the elementary uniform bound $1/(\pi^2a)$ below.
\end{remark}

The exact formula is the classical sine-process number variance
\cite{DysonMehta1963,Mehta2004}; see also
\cite[Theorem~4.3]{Azimifard2026a}. The lemma records its normalization
and explicit remainder for use below. The logarithmic coefficient is
consistent with Landau--Widom~\cite{LandauWidom1980}, and the lattice
analogue with $O(L^{-2})$ error is in
\cite[eq.~(D.6)]{AbanovIvanovQian2011}.

\begin{lemma}\label{lem:two-integrals}
One has
\[
 \int_0^\infty\frac{\sin(2w)\sin^2w}{w^2}\,dw=\log2,
 \qquad
 \int_0^\infty\frac{\cos(2w)\sin^2w}{w^2}\,dw=0,
\]
and for $X>0$ each tail $\int_X^\infty$ is at most $1/X$ in absolute value.
\end{lemma}

\begin{proof}
Since $\sin(2w)\sin^2w=\tfrac12\sin2w-\tfrac14\sin4w=O(w^3)$ at $0$, one
integration by parts gives
$\int_0^\infty(\tfrac12\sin2w-\tfrac14\sin4w)w^{-2}\,dw
=\int_0^\infty(\cos2w-\cos4w)w^{-1}\,dw=\log2$ by Frullani's formula.
Likewise $\cos(2w)\sin^2w=\tfrac14(2\cos2w-1-\cos4w)=O(w^2)$ at $0$, and
$\int_0^\infty\tfrac14(2\cos2w-1-\cos4w)w^{-2}\,dw
=\int_0^\infty(-\sin2w+\sin4w)w^{-1}\,dw=-\tfrac\pi2+\tfrac\pi2=0$.
The tails are bounded by $\int_X^\infty w^{-2}\,dw$.
\end{proof}

\begin{theorem}[Hilbert--Schmidt norm of the critical kernel]\label{thm:hs-exact}
For every $\tau>0$,
\begin{equation}\label{eq:hs-exact}
 \operatorname{tr}\mathcal K_\tau^2
 =4\tau-\frac{2}{\pi^2}\Bigl[\log(4\pi\tau)+1+\gamma+(\log2)\cos(4\pi\tau)\Bigr]
 +R_1(\tau),\qquad |R_1(\tau)|\le\frac{4}{\pi^3\tau}.
\end{equation}
Equivalently $D_J=2\operatorname{tr}(C^{(\pi\tau)}-(C^{(\pi\tau)})^2)
+2\pi^{-2}(\log2)\cos(4\pi\tau)+O(\tau^{-1})$, and
\begin{equation}\label{eq:DJ-DR}
 D_J-2D_R=-\frac{2\log2}{\pi^2}\bigl(1-\cos4\pi\tau\bigr)+O(\tau^{-1}).
\end{equation}
\end{theorem}

\begin{proof}
By the product-to-sum formula,
$\mathcal K_\tau(x,y)=2\cos(\pi\tau(x+y))\,\sin(\pi\tau(x-y))/(\pi(x-y))$,
so with $u=x-y$, $v=x+y$ and $s(u)=\sin^2(\pi\tau u)/(\pi^2u^2)=k_{\pi\tau}(u)^2$,
\[
 \mathcal K_\tau(x,y)^2=4\cos^2(\pi\tau v)\,s(u).
\]
The map $(x,y)\mapsto(u,v)$ satisfies $dx\,dy=\tfrac12\,du\,dv$ and carries $[-1,1]^2$
onto $\{|u|\le2,\ |v|\le L\}$, $L=2-|u|$.  Since
$\int_{-L}^L4\cos^2(\pi\tau v)\,dv=4L+2\sin(2\pi\tau L)/(\pi\tau)$,
\begin{equation}\label{eq:hs-split}
 \operatorname{tr}\mathcal K_\tau^2
 =2\int_{-2}^2(2-|u|)\,s(u)\,du
 +\frac1{\pi\tau}\int_{-2}^2\sin\bigl(2\pi\tau(2-|u|)\bigr)\,s(u)\,du .
\end{equation}
The first term is $2\operatorname{tr}(C^{(\pi\tau)})^2
=2\bigl(2\tau-\operatorname{tr}(C^{(\pi\tau)}-(C^{(\pi\tau)})^2)\bigr)$.
In the second, expand
$\sin(4\pi\tau-2\pi\tau u)=\sin(4\pi\tau)\cos(2\pi\tau u)-\cos(4\pi\tau)\sin(2\pi\tau u)$
and substitute $w=\pi\tau u$:
\begin{align*}
 \int_0^2\sin(2\pi\tau u)\,s(u)\,du
 &=\frac\tau\pi\int_0^{2\pi\tau}\frac{\sin(2w)\sin^2w}{w^2}\,dw
 =\frac\tau\pi\bigl(\log2-\rho_s\bigr),\\
 \int_0^2\cos(2\pi\tau u)\,s(u)\,du&=-\frac\tau\pi\,\rho_c ,
\end{align*}
with $|\rho_s|,|\rho_c|\le(2\pi\tau)^{-1}$ by Lemma~\ref{lem:two-integrals}.
Hence the second term of \eqref{eq:hs-split} equals
\[
 -\frac{2\log2}{\pi^2}\cos(4\pi\tau)
 +\frac{2}{\pi^2}\bigl(\rho_s\cos4\pi\tau-\rho_c\sin4\pi\tau\bigr),
\]
whose remainder is at most $2/(\pi^3\tau)$.  Lemma~\ref{lem:defect-exact}
with $a=\pi\tau$ contributes a further remainder of at most $2/(\pi^3\tau)$,
and \eqref{eq:hs-exact} follows.  For \eqref{eq:DJ-DR} subtract twice
\eqref{eq:defect-asymp} at $a=2\pi\tau$ and use
$\log(4\pi\tau)-\log(8\pi\tau)=-\log2$.
\end{proof}


\subsection{The reflected second moment}
For $0<\ell<1$, write $C_N(\ell)$ for the matrix with entries $c_\ell(m-n)=\sin(\pi\ell(m-n))/[\pi(m-n)]$ and diagonal $\ell$.
\begin{proposition}[Parity-resolved second moment]\label{prop:parity-hs}
For $\tau\ge1$,
\begin{equation}\label{eq:parity-defect}
 \Bigl|\tr(\mathcal R\mathcal K_\tau^2)-\tfrac12\Bigr|
 \le\frac{4\log(2\pi\tau)+2+\pi}{\pi^3\tau}.
\end{equation}
Consequently, by Theorem~\ref{thm:hs-exact}, for $p=\pm1$,
\begin{equation}\label{eq:parity-sector}
\begin{split}
 \tr\bigl(\mathcal K_\tau^2|_{\mathcal H_p}\bigr)
 &=2\tau-\frac1{\pi^2}
    \bigl[\log(4\pi\tau)+1+\gamma+(\log2)\cos(4\pi\tau)\bigr]\\
 &\qquad+\frac p4+O\!\left(\frac{\log\tau}\tau\right).
\end{split}
\end{equation}
The two sectors carry the same bulk $2\tau$ and the same logarithmic
coefficient $-\pi^{-2}$, and separate in the limit only through the constant
$p/4$.  In the discrete frame the reflected second moment agrees with that of the
prolate matrix $C_N(2\varepsilon)$:
\begin{equation}\label{eq:parity-finite}
 \tr\bigl(\mathcal R_NA_N(\varepsilon)^2\bigr)
 =\tr\bigl(\mathcal R_NC_N(2\varepsilon)^2\bigr)
 =\sum_{|m|,|n|\le N}c_{2\varepsilon}(m-n)\,c_{2\varepsilon}(m+n).
\end{equation}
\end{proposition}
\begin{proof}
The product-to-sum identity $2\cos(\pi\tau u)k_{\pi\tau}(u)=k_{2\pi\tau}(u)$
gives the pointwise identity
\[
 K_\tau(-x,y)K_\tau(y,x)
 =k_{2\pi\tau}(x+y)k_{2\pi\tau}(x-y).
\]
The trace of the product of the two Hilbert--Schmidt kernels is their
diagonal integral. With $u=x-y$, $v=x+y$ and Jacobian $1/2$,
\[
 \tr(\mathcal R\mathcal K_\tau^2)
 =\tfrac12\iint_{|u|+|v|\le2}k_{2\pi\tau}(u)k_{2\pi\tau}(v)\,du\,dv.
\]
Integrate in $v$ and substitute $s=2\pi\tau u$. For $T=4\pi\tau$ this gives
\begin{equation}\label{eq:parity-G}
 \tr(\mathcal R\mathcal K_\tau^2)=\frac2{\pi^2}G(T),
 \qquad G(T)=\int_0^T\frac{\sin s}{s}\,\Si(T-s)\,ds.
\end{equation}
Write $G(T)-\tfrac{\pi^2}4
=\int_0^T\frac{\sin s}{s}[\Si(T-s)-\tfrac\pi2]\,ds
 +\tfrac\pi2[\Si(T)-\tfrac\pi2]$ and use
$|\Si(x)-\pi/2|\le\min\{\pi/2,2/x\}$.  The last term is at most $\pi/T$.  In
the integral, bound $T-s\ge T/2$ on $[0,T/2]$ and use
$\int_0^X|\sin s|s^{-1}ds\le1+\log X$ for $X\ge1$; bound $s^{-1}\le2/T$ on
$[T/2,T]$ and split that range at $T-1$, using $2/(T-s)$ on $[T/2,T-1]$ and
$\pi/2$ on $[T-1,T]$.  The three pieces contribute at most
$4(1+\log(T/2))/T$, $4\log(T/2)/T$ and $\pi/T$.  Adding and multiplying by
$2/\pi^2$ gives \eqref{eq:parity-defect}, since $T=4\pi\tau\ge2$.
Then \eqref{eq:parity-sector} follows from
$\tr(\mathcal K_\tau^2|_{\mathcal H_p})
 =\tfrac12[\tr\mathcal K_\tau^2+p\tr(\mathcal R\mathcal K_\tau^2)]$ and
\eqref{eq:hs-exact}.

The identical pointwise calculation on integer nodes, with $\tau$ replaced
by $\varepsilon$, gives
$A_{-m,n}A_{nm}=c_{2\varepsilon}(m+n)c_{2\varepsilon}(m-n)
=C_{-m,n}C_{nm}$, including the removable diagonal cases.
Summation proves \eqref{eq:parity-finite}.
\end{proof}

\begin{corollary}[Leading oscillation of the parity imbalance]
\label{cor:parity-oscillation}
For $\tau\ge1$,
\begin{equation}\label{eq:parity-oscillation}
\begin{split}
 \tr(\mathcal R\mathcal K_\tau^2)
 &=\frac12-
 \frac{[\log(8\pi\tau)+\gamma]\sin(4\pi\tau)}{2\pi^3\tau}
 -\frac{\cos(4\pi\tau)}{4\pi^2\tau}+E_{\rm par}(\tau),\\
 |E_{\rm par}(\tau)|
 &\le\frac{2\log(8\pi\tau)+2\gamma+\pi-5/2}{8\pi^4\tau^2}.
\end{split}
\end{equation}
In particular, the order $O(\log\tau/\tau)$ in
\eqref{eq:parity-defect} is sharp along suitable noninteger scales.
\end{corollary}
\begin{proof}
Put $T=4\pi\tau$, $a(T)=\log(2T)+\gamma$, and $B=\pi/2$.
Differentiating the exact convolution in \eqref{eq:parity-G} gives
$G'(T)=\int_0^T\sin s\sin(T-s)/[s(T-s)]\,ds$.
Symmetry, the identity
$[s(T-s)]^{-1}=T^{-1}[s^{-1}+(T-s)^{-1}]$, and product-to-sum yield
\[
 G'(T)=\frac{\cos T[\operatorname{Ci}(2T)-\gamma-\log(2T)]
                  +\sin T\,\Si(2T)}{T},
 \qquad
 \operatorname{Ci}(y)=-\int_y^\infty\frac{\cos t}{t}\,dt.
\]
Here $\operatorname{Ci}(y)-\gamma-\log y
=\int_0^y(\cos t-1)t^{-1}\,dt$.
Two integrations by parts give, for $y>0$,
\[
\begin{split}
 \operatorname{Ci}(y)
 &=\frac{\sin y}{y}-\frac{\cos y}{y^2}
                   +2\int_y^\infty\frac{\cos t}{t^3}\,dt,\\
 \Si(y)-B
 &=-\frac{\cos y}{y}-\frac{\sin y}{y^2}
                   +2\int_y^\infty\frac{\sin t}{t^3}\,dt.
\end{split}
\]
Thus the errors after their first displayed terms are at most $2/y^2$.
Substitution, with
$\cos T\sin(2T)-\sin T\cos(2T)=\sin T$, gives
\[
 G'(T)=\frac{-a(T)\cos T+B\sin T}{T}
              +\frac{\sin T}{2T^2}+e(T),
 \qquad |e(T)|\le T^{-3}.
\]
Set $F(T)=-[a(T)\sin T+B\cos T]/T$ and
$H(T)=G(T)-\pi^2/4-F(T)$.  Direct differentiation gives
\[
 H'(T)=\frac{[3/2-a(T)]\sin T-B\cos T}{T^2}+e(T).
\]
For $T\ge3$, $b(T)=[a(T)-3/2]/T^2$ is positive and decreasing to zero:
$b'(T)=[4-2a(T)]/T^3<0$.  For any positive decreasing differentiable
weight $b$ tending to zero, integration by parts gives
$|\int_T^\infty b(t)e^{it}\,dt|\le2b(T)$.
Apply this to $b$ and to $B/T^2$; also
$\int_T^\infty|e(t)|\,dt\le1/(2T^2)$.
The already proved estimate \eqref{eq:parity-defect} implies
$G(T)\to\pi^2/4$, while $F(T)\to0$; hence $H(T)\to0$.
Integrating $H'$ from $T$ to infinity therefore proves
\[
 |H(T)|\le\frac{2a(T)-3+\pi+1/2}{T^2}.
\]
Multiply by $2/\pi^2$ in \eqref{eq:parity-G} and substitute
$T=4\pi\tau$ to obtain \eqref{eq:parity-oscillation}.
Along $\tau=n/2+1/8$, the sine equals one and the cosine vanishes,
so the difference from $1/2$ is asymptotic to
$-\log\tau/(2\pi^3\tau)$, proving sharpness of the order.
\end{proof}

